\section{Subarray Distinct Values}
\label{sec:cses-subarray-distinct}

\problemstatement{Given an array of $n$ integers and an integer $k$, count
the number of contiguous subarrays that contain at most $k$ distinct
values.}

\begin{patternbox}
``At most $k$ distinct'' is monotone -- shrinking a window never increases
its distinct count -- so a \textbf{variable-size two-pointer window} works:
for each right end, advance the left end until at most $k$ distinct remain,
and every subarray ending here within the window is valid.
\end{patternbox}

\subsection*{Count subarrays ending at each right end}

Maintain a window $[l,r]$ with a frequency map and its distinct count.
Extend $r$; if adding $a[r]$ pushes the distinct count above $k$, advance
$l$ (removing elements) until it drops back to $k$. Now every subarray that
ends at $r$ and starts anywhere in $[l,r]$ has at most $k$ distinct values,
contributing $r-l+1$ subarrays. Summing that over all $r$ gives the total.

\begin{ideabox}[Window Width Is the Count of Valid Subarrays]
Because the ``at most $k$ distinct'' property is preserved under shrinking,
once the window $[l,r]$ is valid, so is every subarray ending at $r$ with a
start $\ge l$. Adding $r-l+1$ per right end counts them all without
enumeration.
\end{ideabox}

\subsection*{Algorithm}

\begin{enumerate}
  \item $l=0$, empty frequency map, $\textit{distinct}=0$,
        $\textit{ans}=0$.
  \item For $r=0\ldots n-1$: add $a[r]$; while $\textit{distinct}>k$
        remove $a[l]$ and advance $l$; add $r-l+1$ to $\textit{ans}$.
  \item Output $\textit{ans}$.
\end{enumerate}

\subsection*{Dry run}

\dryrun{$a=[1,2,1,3]$, $k=2$.}

\begin{itemize}
  \item $r=0$: window $[1]$, add $1$. $r=1$: $[1,2]$, add $2$. $r=2$:
        $[1,2,1]$ (distinct $2$), add $3$.
  \item $r=3$: adding $3$ gives distinct $3$; shrink to $[2,1,3]$? still
        $3$; shrink to $[1,3]$, distinct $2$; add $2$.
  \item Total $=1+2+3+2=8$ subarrays with at most $2$ distinct.
\end{itemize}

\subsection*{C++ implementation}

\begin{lstlisting}
#include <bits/stdc++.h>
using namespace std;

int main() {
    int n, k;
    cin >> n >> k;
    vector<int> a(n);
    for (int& x : a) cin >> x;

    unordered_map<int,int> freq;
    long long ans = 0;
    int l = 0, distinct = 0;
    for (int r = 0; r < n; r++) {
        if (freq[a[r]]++ == 0) distinct++;
        while (distinct > k) {
            if (--freq[a[l]] == 0) distinct--;
            l++;
        }
        ans += r - l + 1;                      // subarrays ending at r
    }
    cout << ans << "\n";
    return 0;
}
\end{lstlisting}

\begin{complexitybox}
\textbf{Time Complexity.} $O(n)$ amortised -- each pointer moves forward at
most $n$ times. \textbf{Space Complexity.} $O(k)$ distinct keys.
\end{complexitybox}

\begin{takeawaybox}
``Count subarrays with at most $k$ (property)'' is a variable window where
each right end contributes its window width. To count \emph{exactly} $k$
distinct, take atMost$(k)-$atMost$(k-1)$ -- a standard follow-up on the same
routine.
\end{takeawaybox}
